%{{{PREAMBLE
\documentclass[12pt,b5paper]{article}
\usepackage[utf8]{inputenc}
\usepackage{amsmath}
\usepackage{amsfonts}
\usepackage{amssymb}
\usepackage{xparse}
\usepackage{comment}
\usepackage{tikz}
\usepackage{mathrsfs}



\begin{comment}
\usepackage[margin=0.70in]{geometry}
\usepackage[utf8]{inputenc}
\usepackage[T1]{fontenc}
\usepackage[english]{babel}
\usepackage{graphicx, float, subfigure, blindtext}
\usepackage[\IEEEhyperrefsetup, pdftex]{hyperref}
\usepackage{mathtools}
\usepackage{amsthm}
\usepackage{bm}
\usepackage{amsmath, mathdesign}
\usepackage{amsfonts}
\usepackage{amssymb}
\usepackage{geometry}
\usepackage{mathrsfs}
\usepackage{yfonts}
\usepackage{forest}
\usepackage{cuted}
\usepackage[nolist,nohyperlinks]{acronym}
\usepackage{cleveref}
\usepackage{titlesec}
\end{comment}
\author{Adrian Swande}
\begin{document}
%}}}
%{{{CUSTOM COMMANDS
%\newcommand{\vcat}{+\kern-5px\circ\ } %\text{‡}} %\bbox[black]{\color{white}{\bullet}} }
%\newcommand{\vcat}{
%  \mathord{
%    \hbox{
%      \ooalign{
%		\hfil$\displaystyle +$\hfil\cr
%		\hfil\raisebox{0ex}{\smash{
%			%$\times$
%		\ooalign{
%		\hfil$\displaystyle \times$\hfil\cr
%		\hfil\raisebox{0ex}{\smash{$\bullet$}{}}\hfil\cr
%      }
%}{}}\hfil\cr
%      }
%    }
%  }
%}
\newcommand{\vcat}{
  \mathord{
    \hbox{
      \ooalign{
		  \hfil$\displaystyle {\circ}$\hfil\cr
		\hfil\raisebox{0ex}{\smash{$*$}{}}\hfil\cr
      }
    }
  }
}
\newcommand{\quotemath}[1]{\text{``$#1$''}}
%\newcommand{\vect}[1]{\left\langle #1 \right\rangle}
\newcommand{\vect}[1]{\left[ #1 \right]}
\newcommand{\svect}[1]{\begin{pmatrix} #1 \end{pmatrix}}
\newcommand{\set}[1]{\left\{ #1 \right\}}
\newcommand{\setdef}[2]{\set{#1 \,\middle|\, #2}}
\newcommand\bbigcirc{\scalebox{1.5}{$\bigcirc$}}
\newcommand{\condw}[1]{\left\{\begin{array}{cl} #1 \end{array}\right.}
\newcommand{\conde}[2]{#1 & : \ #2 \\}
\newcommand{\prnth}[1]{\left(#1\right)}
\newcommand{\solarr}{\mathbf{s}}
\newcommand{\solperm}{\acute{\solarr}}
\newcommand\bbigcap{\scalebox{1.5}{$\bigcap$}}
\newcommand\bbigcup{\scalebox{1.5}{$\bigcup$}}
%}}}

\section{Data Generation}

\subsection{Pattern Matching}

A \textit{pattern} is a set.

%A \textit{pattern of }

\subsubsection{Example}

Let \(a \vcat b= \vect{a_1,\cdots,a_n,b_1,\cdots,b_m}\), where \(a=\vect{a_1,\cdots,a_n}\) and \(b=\vect{b_1,\cdots,b_m}\).

\[
\begin{array}{ccl}
	\mathfrak{E} & := &
	\mathfrak{T} \cup \setdef{\mathfrak{e}_1 \vcat \mathfrak{o} \vcat \mathfrak{e}_2}{\mathfrak{o}\in\set{\vect{\quotemath{+}},\vect{\quotemath{-}}} \land \mathfrak{e}_1,\mathfrak{e}_2 \in \mathfrak{E}} \\

	\mathfrak{T} & := &
	\mathfrak{F} \cup \setdef{\mathfrak{t}_1 \vcat \mathfrak{o} \vcat \mathfrak{t}_2}{\mathfrak{o}\in\set{\vect{\quotemath{\times}},\vect{\quotemath{\div}}} \land \mathfrak{t}_1,\mathfrak{t}_2 \in \mathfrak{T}} \\

	\mathfrak{F} & := &
	\mathfrak{N} \cup \setdef{\vect{\quotemath{(}} \vcat \mathfrak{e} \vcat \vect{\quotemath{)}}}{\mathfrak{e} \in \mathfrak{E}} \\

	\mathfrak{N} & := &
	N \cup N^* \\

	N & := & \set{\text{\quotemath{0}}, \cdots, \text{\quotemath{9}}}
\end{array}
\]

\textit{What is the difference between pattern matching and tokenwize parsing?}

\newpage
\section{\texttt{gen.jl} and \texttt{pattern.jl}}

The \textit{map form} of a pattern \(\mathfrak{P}\) is a function \(\Phi':\mathfrak{P}\to L\), where \(L\subseteq\mathcal{P}(\mathfrak{P})\) is a set of \textit{labels}\footnote{Lables -- or, \textit{concepts} -- are useful abstractions of patterns, which will be used later to optimize the complexity of solution-finding systems. The subsets \(\mathfrak{T, F, N}\) and \(N\) in example 1.1.1 are examples of lables.}.
%\textit{concepts}.
The \textit{inverse map form} of a pattern \(\mathfrak{P}\) is a function \(\Phi:L \to \mathcal{P}(\mathfrak{P})\), where \(\Phi(l) = \setdef{p\in\mathfrak{P}}{\Phi'(p)=l}\) for \(l\in L\).

Let \(G:=\left( \left\langle \mathscr{G} \right\rangle, * \right)\) be a group. Let \(\triangle:\Phi \times L^2 \to \Phi\) \color{red}!!!!!!!!!!!!WRONG, \(\Phi\) should be replaced by a \textit{set} of \(\Phi\)'s!!!!!!!!!!!\color{black} -- where \(\Phi\) is the inverse map form of a pattern \(\mathfrak{P} \subseteq G^{<\omega}\) and \(L = \text{dom}(\Phi) \subseteq G\) -- be defined by 
\(\forall l\in\text{dom}(\Phi_1)\setminus\set{l_1*l_2}\ \left(\Phi_2(l) = \Phi_1(l)\right)\) and \(\Phi_2(l_1*l_2) = \Phi_1(l_1*l_2)\cup\set{\vect{l_1,l_2}}\) if \(l_1*l_2\in\text{dom}(\Phi_1)\) and \(\Phi_2(l_1*l_2) = \set{\vect{l_1,l_2}}\) if not -- where \(\Phi_2 = \triangle(\Phi_1, \vect{l_1, l_2})\)\color{red}????????\color{black} and \(l_1, l_2 \in \text{dom}(\Phi_1)\).
%\(\forall l\in\text{dom}(\Phi_1)\ \left(\Phi_1(l) \subseteq \Phi_2(l)\right)\) and \(\vect{l_1,l_2} \in \Phi_2(l_1*l_2)\) where \(\Phi_2 = \triangle(\Phi_1, \vect{l_1, l_2})\) and \(l_1, l_2 \in \text{dom}(\Phi_1)\).

Now, consider the function \(\mathscr{R}:\Phi \times \mathbb{N} \to \Phi\) \color{red}!!!!!!!!!\color{black} defined by 
\[
	\mathscr{R}(\Phi, n) = \condw{
	\conde{\Phi}{n=0}
	\conde{\mathscr{R}\prnth{\triangle\prnth{\Phi, \overset{2}{\underset{i=1}{\scalebox{2}{$\vcat$}}}\text{Rand}\prnth{\text{dom}\prnth{\Phi}}}}, n-1)}{\text{otherwise}}
}.
\]
Further, \(\Phi_n = \mathscr{R}(\Phi, n)\), where 
%\(\text{dom}(\Phi) = \mathscr{G}\) and \(\text{Img}(\Phi) = \set{\emptyset}\) 
\(\Phi:\mathscr{G} \to \set{\emptyset}\) 
and \(n\in\mathbb{N}\), is a partial or full representation of the group \(G\), where each element in the domain is mapped thereby to the set of its neighbours.

For each \(l\in\text{dom}(\Phi_n)\), then, a set of vectors \(\solarr\in\mathscr{G}^{<\omega}\) can be inferred with \(\kappa:\Phi\times G\to\mathscr{G}^{<\omega}\)\color{red}!!!!!!!!! + should \(\mathscr{G}\) be included in domain too? probably, yes!!!!!!!!\color{black}, defined by 
\[
	\kappa(\Phi, g) = \condw{
		\conde{g}{g\in\mathscr{G}}
		%\conde{\underset{k\in\kappa\prnth{\Phi, \Phi\prnth{o}}}\bigcup k}{\text{otherwise}}
		\conde{\underset{k\in\Phi(g)}\bigcup {\prnth{\underset{j\in k}{\scalebox{2}{$\vcat$}}\kappa\prnth{\Phi, j}}} }{\text{otherwise}}
	}
\]
\color{red}!!!!!!!!!!!THE DEFINITION IS WRONG, NOW THE VECTORS WILL INCLUDE SETS OF POSSIBLE VECTORS!!!!!!!!!!!\color{black}

\end{document}


% TODO
% - figure out how to insert \koppa, and replace \kappa therewith so as to match the julia programs (gen.jl)
% - model weights
% - Meta: treat \Phi's as elements in a group with \triangle ( + [l_1,l_2]) as an operator. Then IMF patterns can be constructed using the same system -- IMF patterns for adding stuff to IMF patterns. What this is, is observing the change in patterns from multiple puzzles/worlds. -- this can/should be expanded further, with observing observations of change in exploring puzzles and so on...
